\documentclass[letterpaper,twocolumn,10pt]{article}
\usepackage{usenix,epsfig,endnotes,amsmath}
\usepackage{multirow,booktabs,graphicx,subcaption}
\begin{document}

\date{}

\title{\Large \bf Detection of Lateral Movement in Cloud VPC Networks via Graph-Based Analysis of Flow and Authentication Logs}

\author{
{\rm Ibrahim Pehlivan}\\
University of California, Los Angeles
\and
{\rm Wesley Gunawan}\\
University of California, Los Angeles
\and
{\rm Samyak Kakatur}\\
University of California, Los Angeles
}

\maketitle

\thispagestyle{empty}


\subsection*{Abstract}
Lateral movement in cloud VPCs is difficult to detect because internal traffic is implicitly trusted and individual log sources provide incomplete visibility. We present a graph-based detection method that combines network flow logs and authentication logs into a unified directed multigraph, extracting structural, temporal, and statistical features to score edges for lateral movement likelihood. We evaluate three data variants (auth-only, combined, flow-only) with three detection approaches (graph-based scoring, Isolation Forest, One-Class SVM) on the LANL-2015 dataset. Our graph-based auth-only method achieves 92.2\% recall at 0.64\% false positive rate (FPR) with 0.984 AUC, outperforming Isolation Forest (0.973 AUC) and One-Class SVM (0.923 AUC) on the same features. Graph-based scoring maintains AUC $\geq$ 0.96 across all data variants, while unsupervised baselines degrade to 0.66--0.72 AUC on flow-only data, showing that graph structure carries signal that tabular features alone lose. Graph-derived features alone achieve up to 0.998 AUC compared to 0.823 for pure tabular features on flow data, a 21\% gain. The pipeline processes 23,000--73,000 events per second.

\section{Introduction}

\begin{figure}[htbp]
\centering

\includegraphics[width=0.95\columnwidth]{lat_mvmt.png}
\caption{Modified Advanced-Persistent-Threat Model, showing Lateral Movement as the step after the attacker compromises one machine (Foothold). Adapted from Mamun et al.~\cite{APT-detection}.}
\label{fig:lat_mvmt}
\end{figure}

Lateral movement is the phase of a network attack where an adversary, having gained an initial foothold, pivots through a compromised network to reach high-value targets such as domain controllers, databases, or exfiltration endpoints. In cloud VPC (Virtual Private Cloud) environments, lateral movement is especially difficult to detect because internal traffic flows freely between hosts by default, and the volume of benign connections dwarfs the few attacker-initiated ones.

No single log source provides enough signal to reliably identify lateral movement. Network flow logs capture connections but cannot distinguish an attacker's session from an administrator's. Authentication logs record login events but are blind to reconnaissance activity such as port scans or service enumeration that never triggers a login event. Existing perimeter-focused intrusion detection systems are designed to identify north-south traffic crossing the network boundary, not the east-west movement that characterizes lateral movement within a cloud VPC.

We address this gap by constructing a unified streaming graph from both flow and authentication logs, where computers and user identities are nodes, and connections and authentication events are edges. Graph-based features (degree centrality, fan-out ratios, temporal patterns, and betweenness) capture the relational structure of lateral movement that tabular features miss. Our detection pipeline scores each edge for lateral movement likelihood, applies path-level analysis to identify multi-hop attack chains, and thresholds scores to flag anomalous behavior. On auth-only data this configuration reaches 92.2\% recall at 0.64\% FPR, and Figure~\ref{fig:roc_all} shows graph-based scoring holding above 0.96 AUC on every variant while unsupervised baselines drop on flow-only data (Table~\ref{tab:baseline_comparison}).

The main contributions of this paper are:

\begin{enumerate}
\item \textbf{Unified graph framework:} we integrate flow and authentication logs into a single directed multigraph for lateral movement analysis.
\item \textbf{Multi-variant evaluation:} graph-based auth-only analysis reaches 92.2\% recall at 0.64\% FPR (Table~\ref{tab:methods_comparison}), beating Isolation Forest and One-Class SVM on the same features (Table~\ref{tab:baseline_comparison}).
\item \textbf{Feature ablation:} graph-derived features yield a 21\% AUC gain over tabular features on flow data (Section~\ref{sec:ablation}).
\item \textbf{False positive analysis:} we identify the structural causes of false alarms and propose targeted mitigation strategies (Section~\ref{sec:fp_analysis}).
\end{enumerate}

Across the full evaluation, graph-based auth-only detection reaches 0.984 AUC at 0.64\% FPR (Table~\ref{tab:methods_comparison}), and Figure~\ref{fig:roc_all} plots ROC curves for every variant. Section~\ref{sec:ablation} traces most of the discriminative signal to graph topology rather than tabular features.

\section{Related Work}

Graph-based detection methods model relationships between system events and reason over paths or substructures. Hopper~\cite{hopper2021} builds a login graph from authentication events and uses path inference to flag lateral movement sequences. However, its reliance on authentication logs alone means it cannot detect lateral movement that bypasses login entirely, such as exploitation of services or use of stolen session tokens. We extend this idea by incorporating network flow logs, capturing lateral movement paths that never produce an authentication event. Euler~\cite{euler2023} applies temporal link prediction to cloud VPC network graphs, treating lateral movement as a link prediction problem in an evolving graph. Its temporal modeling is a strong baseline for our own temporal graph features, though it operates on a single data source. POIROT~\cite{poirot2019} matches provenance graph templates against audit logs to detect known attack patterns. Template matching is effective when attack structures are known in advance but cannot catch novel variations; we use anomaly-based scoring instead and adapt the graph model to cloud VPC network logs rather than host-level audit data.

Multi-log correlation approaches combine evidence from separate data sources to improve detection coverage. HOLMES~\cite{holmes2019} correlates system audit logs across multiple stages of the kill chain, mapping observed behavior to MITRE ATT\&CK tactics and producing a real-time threat score. Its multi-stage correlation confirms that combining log sources can improve detection, and its kill-chain framing shaped our detection pipeline structure. However, HOLMES operates on host-level audit logs in traditional enterprise settings, whereas we target network flow logs and authentication events at cloud VPC network scale.

\section{Methodology}

We evaluate whether a graph built from both network flow and authentication logs detects lateral movement more effectively than graphs from either source alone, and whether graph-based scoring outperforms standard unsupervised anomaly detectors. We build three graph variants per run: \emph{auth-only}, \emph{flow-only}, and \emph{combined}, each scored by graph-based weighting, Isolation Forest, and One-Class SVM.

\subsection{Graph Construction}

We model the network as a directed multigraph. Nodes represent computers and users; edges represent authentication events or network flows. A streaming builder reads gzipped CSV files line-by-line, keeping only events within time windows centered on red-team activity ($\pm$3,600 seconds per event, merged into 25 non-overlapping intervals). Events outside all windows are skipped, so memory stays proportional to active windows, not the full 1.6 billion-event dataset.

For each authentication event, the builder adds a computer-to-computer edge storing auth type, logon type, orientation, success, and timestamp. For each flow event, a computer-to-computer edge stores protocol, source/destination ports, packet count, byte count, duration, and timestamp. When multiple events share the same $(src, dst)$ pair, the edge weight increments and only the first and last timestamps are kept. Node types are inferred by naming convention: names containing \texttt{\$} before \texttt{@} or ending with \texttt{\$} are machine accounts; everything else is a user.

\subsection{Feature Extraction}

We extract 23 features across three granularities: per-edge (11), per-node (7), and per-graph (5).

\textbf{Edge features (11).}
Each edge receives: edge rarity ($1 / \text{weight}$, percentile-ranked), source out-degree, destination in-degree, destination fan-out ratio, source fan-out, normalized weight, self-loop flag, user-edge flag, and a temporal decay weight. For authentication edges: whether the auth type is NTLM (\texttt{is\_ntlm}), whether the logon type is Network (\texttt{is\_network\_logon}), and whether authentication succeeded (\texttt{is\_success\_auth}). For flow edges: whether the destination port is unusual (\texttt{is\_unusual\_dst\_port}); protocol rarity; byte-per-packet ratio (percentile-ranked); and duration z-score.

\textbf{Node features (7).}
For each node we compute in-degree, out-degree, total degree, and fan-out ratio (out-degree / total degree). The fan-out ratio captures a known lateral movement pattern: compromised hosts performing reconnaissance connect to many distinct destinations (high fan-out), while normal servers mostly receive connections (low fan-out). Betweenness centrality is computed with an approximation cutoff for large graphs. We also compute four temporal features per node from its outgoing edges: mean and standard deviation of inter-arrival times, a burst score (maximum fraction of outgoing edges falling within any 10\% sub-window of the node's active span), and total active duration. Node degree features are projected onto edges so that flat baselines can use connectivity information.

\textbf{Graph features (5).}
Density, average clustering coefficient (computed on the undirected version), weakly connected component count, node count, and edge count.

\subsection{Feature Selection and Auditing}

To select the feature subset for each graph variant, we audit each candidate feature with a held-out AUC protocol. The cached edge table is joined with node-level statistics, filtered to remove self-loops and user-edge metadata, and labeled by whether each source-destination pair appears in the LANL red-team ground truth. A stratified 50/50 split reserves one half for calibrating the feature ranking and the other half for checking generalization. For each feature, we compute ROC AUC using that feature alone; if red-team edges have a lower mean than benign edges, we invert the AUC so the score measures separability regardless of direction.

The top features selected per variant are:

\begin{itemize}
\item \textbf{Auth-only (top 3):} \texttt{src\_in\_degree} (AUC 0.940), \texttt{is\_ntlm} (AUC 0.909), \texttt{src\_out\_degree} (AUC 0.893)
\item \textbf{Combined (top 5):} \texttt{is\_ntlm} (AUC 0.933), \texttt{dst\_in\_degree} (AUC 0.819), \texttt{is\_network\_logon} (AUC 0.817), \texttt{dst\_total\_degree} (AUC 0.812), \texttt{edge\_rarity} (AUC 0.808)
\item \textbf{Flow-only (top 3):} \texttt{src\_burst\_score} (AUC 0.807), \texttt{dst\_inter\_arrival\_mean} (AUC 0.799), \texttt{dst\_inter\_arrival\_std} (AUC 0.777)
\end{itemize}

\begin{figure}[htbp]
\centering
\includegraphics[width=0.95\columnwidth]{feature_audit_combined.png}
\caption{Feature audit for the combined variant. NTLM authentication is the strongest individual signal (AUC 0.93), followed by destination in-degree and network logon type (both 0.82). Seven of ten features exceed 0.75 AUC, indicating strong individual discriminative power for lateral movement detection.}
\label{fig:feat_audit}
\end{figure}

Figure~\ref{fig:feat_audit} shows the individual AUC scores for all audited features in the combined variant. The \texttt{src\_betweenness\_centrality} feature scores 0.17, indicating it is inversely correlated with lateral movement edges: attackers tend to target peripheral rather than central nodes, consistent with the LANL red-team strategy of compromising a single jump box (C17693) and pivoting from there.

\subsection{Weight Optimization and Held-Out Validation}

We optimize the selected edge features using the Nelder-Mead algorithm to maximize AUC. For each graph variant, the optimizer converges independently on the full dataset, then we validate generalization using a stratified 50/50 calibration/evaluation split (seed 42):

\begin{itemize}
\item \textbf{Auth-only:} AUC improves from 0.907 (equal weights) to 0.985 (+8.6\%). Held-out: calibration 0.956, evaluation 0.976. Optimized weights: \texttt{src\_out\_degree} = 0.809, \texttt{is\_ntlm} = 0.379, \texttt{src\_in\_degree} = $-$0.449.
\item \textbf{Combined:} AUC improves from 0.862 (equal weights) to 0.959 (+11.3\%). Held-out: calibration 0.965, evaluation 0.963. Optimized weights: \texttt{is\_ntlm} = 0.746, \texttt{is\_network\_logon} = 0.627, \texttt{dst\_total\_degree} = $-$0.487, \texttt{edge\_rarity} = 0.028.
\item \textbf{Flow-only:} AUC improves from 0.929 (equal weights) to 0.963 (+3.6\%). Held-out: calibration 0.765, evaluation 0.768.
\end{itemize}

The negative weight on \texttt{src\_in\_degree} (auth-only) and \texttt{dst\_total\_degree} (combined) indicates that lateral movement edges target less-connected hosts rather than high-traffic servers, consistent with attackers avoiding noisy, well-monitored systems.

\begin{figure}[htbp]
\centering

\includegraphics[width=0.95\columnwidth]{holdout_validation.png}
\caption{Holdout validation: Nelder-Mead optimized weights vs. logistic regression on a 50/50 stratified split. Both methods generalize with gaps under 2.1\%. Auth-only achieves the highest held-out AUC (0.976). The near-identical performance confirms that the AUC derives from feature quality, not optimization method.}
\label{fig:holdout}
\end{figure}

Figure~\ref{fig:holdout} shows the held-out validation results. Logistic regression on the same features achieves nearly identical eval AUC: combined 0.962, auth-only 0.976, flow-only 0.758. This confirms that the AUC improvement comes from the features themselves: both Nelder-Mead direct AUC maximization and logistic regression converge to essentially the same performance.

\subsection{Edge Scoring and Path Analysis}

Edges are scored with a weighted sum of the optimized features. All non-binary features are percentile-rank transformed before the weighted sum. Self-loops and user-to-user edges are explicitly zeroed (score = 0.0).

Single-edge scoring misses multi-hop lateral movement chains. To capture these, we enumerate paths via breadth-first search from every node, up to 4 hops, following only the top outgoing edges per node ranked by edge score. Each path is scored as the mean of its edge scores' geometric mean, maximum, and arithmetic mean. The top 50 paths are retained, and every edge appearing in a top-50 path receives a boost of $0.1 \times \text{path\_score}$ added to its own score, capped at 1.0. This feeds path-level lateral movement signal back into edge-level detection.

\subsection{Threshold Selection}

Anomalous edges are identified by thresholding final edge scores. We auto-search over percentile thresholds [90, 95, 97, 99, 99.5, 99.9] and select the one that maximizes F1 against known red-team pairs. The optimal thresholds are: auth-only at the 99th percentile (score = 1.1179), combined at the 99.5th percentile (score = 1.2686), and flow-only at the 97th percentile (score = 0.6952).

\subsection{Baseline Methods}

We compare graph-based scoring against Isolation Forest~\cite{isoforest} and One-Class SVM. Both operate on the same edge feature vector, using scikit-learn defaults (Isolation Forest: n\_estimators=100, contamination="auto"; One-Class SVM: kernel="rbf", nu=0.1). Features are standardized via \texttt{StandardScaler}.

\section{Experiment Setup}

\subsection{Dataset}

\textbf{LANL Cybersecurity Dataset (LANL-2015)~\cite{kent-2015-cyberdata1}.}
The dataset covers 58 days of continuous network activity across 12,425 users and 17,684 computers in three gzipped CSV files: authentication events, network flow events, and 749 labeled red-team attack events forming 308 unique source-destination pairs. Red-team events include lateral movement via SMB, SSH, RDP, credential reuse, and network reconnaissance; 93.6\% originate from a single compromised host (C17693) used as a jump box.

Because loading 1.6 billion events is impractical, we scope the data using time windows. For each red-team event, we define a window of $\pm$3,600 seconds. Overlapping windows are merged, producing 25 non-overlapping intervals that capture approximately 104 million auth events and 31.6 million flow events.

\subsection{Data Characteristics}

Attack traffic in LANL-2015 is nearly all TCP (99.82\% for red-team flows vs. roughly 50\% for benign traffic). The fan-out ratio captures a known lateral movement pattern: compromised hosts performing reconnaissance connect to many distinct destinations (high fan-out), while normal servers mostly receive connections (low fan-out). Attackers exhibit markedly different inter-arrival patterns than normal users. These patterns are relational: they emerge from connections between entities, not from individual events in isolation.

\subsection{Evaluation Metrics}

All detection metrics are computed in pair-space: after thresholding, anomalous edges are collapsed into source-destination pairs, which are compared against the 308 ground-truth red-team pairs. AUC is computed over all valid edges (not pairs) and is threshold-independent: it measures whether scoring ranks lateral movement edges above normal ones regardless of the chosen cutoff.

With 308 red-team pairs among hundreds of thousands of total edges, precision will always be low. The meaningful comparisons are AUC (ranking quality) and FPR at matched recall (operational cost of detection).

\subsection{Implementation}

The pipeline is implemented in Python using igraph for graph operations and scikit-learn for baselines and weight optimization. All hyperparameters are stored in JSON configuration files. The pipeline processes 23,000--73,000 events per second depending on log source, completing end-to-end in roughly 36 minutes with 20 parallel workers.

\section{Results}

\subsection{ROC Curves by Variant}

\begin{figure}[htbp]
\centering
\begin{subfigure}[b]{0.95\columnwidth}
\centering
\includegraphics[width=\textwidth]{roc_curves_auth_only.png}
\caption{Auth Only}
\label{fig:roc_auth}
\end{subfigure}
\begin{subfigure}[b]{0.95\columnwidth}
\centering
\includegraphics[width=\textwidth]{roc_curves_combined.png}
\caption{Combined}
\label{fig:roc_combined}
\end{subfigure}
\begin{subfigure}[b]{0.95\columnwidth}
\centering
\includegraphics[width=\textwidth]{roc_curves_flow_only.png}
\caption{Flow Only}
\label{fig:roc_flow}
\end{subfigure}
\caption{ROC curves for all three methods on each data variant. Graph-based scoring achieves the highest AUC on all three: 0.98 (auth-only), 0.96 (combined), and 0.96 (flow-only). The gap is widest on flow-only data, where unsupervised baselines degrade to 0.66--0.72 AUC.}
\label{fig:roc_all}
\end{figure}

Figure~\ref{fig:roc_all} presents ROC curves for all three methods on each data variant. Graph-based scoring holds above 0.96 AUC on flow-only data while both unsupervised baselines collapse. The reason is structural. Lateral movement sessions look much like benign ones at the per-flow level: the same ports, protocols, and byte counts appear in both. Topology, which encodes which hosts connect to which, does not, and that relational structure is exactly what the baselines discard.

\subsection{Graph-Based Method Performance}

\begin{table}[h]
\centering
\begin{tabular}{lcccccc}
\toprule
Method & Recall & FPR & F1 & AUC & Latency (s) & Throughput (events/s) \\
\midrule
auth\_only  & \textbf{0.9221} & 0.0064 & \textbf{0.1778} & \textbf{0.9842} & 3903.87 & 26,659 \\
flow\_only  & 0.0682 & 0.0298 & 0.0099 & 0.9626 & 430.02 & \textbf{73,511} \\
combined    & 0.2110 & \textbf{0.0036} & 0.0591 & 0.9590 & 5899.24 & 23,000 \\
\bottomrule
\end{tabular}
\caption{Graph-based method comparison. Auth-only achieves the highest recall (92.2\%) and AUC (0.984). Combined achieves the lowest FPR (0.36\%) but at the cost of recall (21.1\%). Flow-only is fastest but captures few lateral movement pairs.}
\label{tab:methods_comparison}
\end{table}

\begin{table}[h]
\centering
\small
\resizebox{\columnwidth}{!}{%
\begin{tabular}{lcccccc}
\toprule
Method & Graph Nodes & Graph Edges & RT Pairs in Graph & Anomalous Pairs & Threshold & Score \\
\midrule
auth\_only  & 90,783 & 409,100 & 302 & 2,886 & 99th & 1.1179 \\
combined    & 91,589 & 512,529 & 305 & 1,892 & 99.5th & 1.2686 \\
flow\_only  & 11,586 & 131,562 & 35 & 3,947 & 97th & 0.6952 \\
\bottomrule
\end{tabular}%
}
\caption{Graph structure and detection statistics. Auth-only and combined capture nearly all red-team pairs (302--305 of 308). Flow-only captures only 35, as many attack endpoints are invisible through flow logs alone.}
\label{tab:graph_stats}
\end{table}

Table~\ref{tab:methods_comparison} shows that auth-only is the strongest variant for lateral movement detection: it captures 302 of 308 red-team pairs and reaches 92.2\% recall at 0.64\% FPR. The combined method reaches the lowest FPR but its score is diluted by flow edges that do not carry strong attack signals, cutting recall to 21.1\%. Flow-only is structurally limited: it captures only 35 of 308 red-team pairs because most attack endpoints never produce a flow-level signature the scoring can exploit.

Table~\ref{tab:graph_stats} reports the underlying graph sizes and pair counts. The flow-only graph is an order of magnitude smaller than the other two and reaches only 35 of the 308 red-team pairs, which is the structural reason its recall cannot match auth-only.

\subsection{Baseline Comparison: Graph vs. Unsupervised Methods}

\begin{figure}[htbp]
\centering
\includegraphics[width=0.95\columnwidth]{method_comparison_auc.png}
\caption{AUC comparison across all methods and variants. Graph-based scoring maintains AUC $\geq$ 0.95 across all variants, while unsupervised baselines degrade significantly on flow-only data.}
\label{fig:method_auc}
\end{figure}

\begin{table}[h]
\centering
\begin{tabular}{lccccc}
\toprule
Variant & Method & AUC & Recall & FPR & F1 \\
\midrule
\multirow{3}{*}{Auth Only}
  & Graph-based      & \textbf{0.9842} & \textbf{0.9221} & \textbf{0.0064} & \textbf{0.1778} \\
  & Isolation Forest & 0.9726 & 0.9448 & 0.0291 & 0.1322 \\
  & One-Class SVM    & 0.9234 & 0.9416 & 0.1068 & 0.0399 \\
\midrule
\multirow{3}{*}{Combined}
  & Graph-based      & \textbf{0.9590} & 0.2110 & \textbf{0.0036} & 0.0591 \\
  & Isolation Forest & 0.9572 & 0.5877 & 0.0291 & \textbf{0.0626} \\
  & One-Class SVM    & 0.8000 & 0.0519 & 0.0009 & 0.0652 \\
\midrule
\multirow{3}{*}{Flow Only}
  & Graph-based      & \textbf{0.9626} & \textbf{0.0682} & \textbf{0.0298} & \textbf{0.0099} \\
  & Isolation Forest & 0.7182 & 0.0422 & 0.0999 & 0.0038 \\
  & One-Class SVM    & 0.6622 & 0.0292 & 0.1018 & 0.0026 \\
\bottomrule
\end{tabular}
\caption{Full comparison of graph-based scoring against unsupervised baselines on the same edge features. Graph-based achieves the highest AUC on all three variants. The gap is most pronounced on flow-only data (0.963 vs. 0.718 for IF, 0.662 for SVM). Graph-based recall and FPR are reported in pair-space (anomalous edges collapsed to source--destination pairs); baseline recall and FPR are reported in edge-space and are not directly comparable.}
\label{tab:baseline_comparison}
\end{table}

Figure~\ref{fig:method_auc} and Table~\ref{tab:baseline_comparison} compare graph-based scoring against two unsupervised baselines on the same edge features.

\begin{itemize}
\item \textbf{On auth-only data}, all three methods find most lateral movement pairs, but graph-based holds the lowest pair-space FPR by a wide margin (Table~\ref{tab:baseline_comparison}). Fewer false alerts per real detection matters operationally, since a SOC analyst triages every flagged pair.

\item \textbf{On combined data}, graph-based and Isolation Forest are nearly tied in AUC, yet graph-based again reports a much lower pair-space FPR while One-Class SVM degrades sharply. For mixed-source deployments, the graph formulation keeps the alert volume manageable.

\item \textbf{On flow-only data}, the gap is largest: graph-based stays above 0.96 AUC while both baselines fall to 0.66--0.72. Lateral movement relies on which hosts reach which, not on per-flow attributes, so detectors that ignore topology lose the signal entirely.
\end{itemize}

\begin{figure}[htbp]
\centering
\includegraphics[width=0.95\columnwidth]{metrics_summary.png}
\caption{Performance metrics summary across all methods and variants. Best value per metric highlighted. Graph-based auth-only achieves the best AUC (0.984) and F1 (0.178).}
\label{fig:metrics_summary}
\end{figure}

\subsection{Graph Topology Differences}

\begin{figure}[htbp]
\centering
\includegraphics[width=0.95\columnwidth]{graph_statistics.png}
\caption{Graph topology metrics across data variants (log scale). Auth-only and combined have similar scale (90K nodes, 400K--500K edges). Flow-only is smaller (11K nodes) but denser, with 2--3$\times$ higher average degree.}
\label{fig:graph_stats}
\end{figure}

Figure~\ref{fig:graph_stats} reveals how the graph structure differs across data variants. The auth-only and combined graphs are similar in scale (90K nodes, 400K--500K edges), reflecting the large number of user-machine authentication relationships in an enterprise network. The flow-only graph is an order of magnitude smaller (11K nodes, 131K edges) but denser, with 2--3$\times$ higher average degree. This structural difference explains why graph features are most valuable for flow data: the smaller, denser flow graph has richer connectivity patterns per node that graph-derived features can exploit, while tabular features alone lack the context to distinguish attack flows from the dense web of normal traffic.

\begin{figure}[htbp]
\centering
\includegraphics[width=0.95\columnwidth]{radar_chart_auth_only.png}
\caption{Radar profile for auth-only variant: all methods achieve high Recall, AUC, and 1-FPR, but diverge significantly on F1. Graph-based achieves the highest F1 (0.178).}
\label{fig:radar_auth}
\end{figure}

Figure~\ref{fig:radar_auth} shows the multi-metric radar profile for the auth-only variant. All three methods achieve high Recall, AUC, and 1-FPR (specificity), clustering near the outer edge. The divergence occurs on F1, where graph-based (0.178) significantly outperforms Isolation Forest (0.132) and One-Class SVM (0.040). This pattern reflects the class imbalance: all methods find most lateral movement pairs (high recall), but graph-based generates fewer false positives, yielding higher precision and thus higher F1.

\subsection{Tabular vs. Graph Feature Ablation}
\label{sec:ablation}

\begin{figure}[htbp]
\centering

\includegraphics[width=0.95\columnwidth]{ablation_study.png}
\caption{Feature ablation: pure tabular vs. graph-derived vs. combined features under logistic regression holdout. Graph-derived features dominate across all three data variants. The largest gap is on flow-only data (0.998 vs. 0.823 AUC).}
\label{fig:ablation}
\end{figure}

Figure~\ref{fig:ablation} presents the ablation study results evaluated under logistic regression with a stratified 50/50 holdout split. Graph-derived features outperform pure tabular features across all three variants, with the most dramatic gap on flow-only data:

\begin{itemize}
\item \textbf{Flow-only:} Graph-derived 0.998 vs. tabular 0.823 (+21.3\%). Graph structure transforms weak flow signals into near-perfect lateral movement detection.
\item \textbf{Auth-only:} Graph-derived 0.991 vs. tabular 0.963 (+2.9\%). Even on the strongest data source, graph features provide a meaningful gain.
\item \textbf{Combined:} Graph-derived 0.987 vs. tabular 0.953 (+3.8\%). Adding tabular features to graph features yields negligible additional benefit (0.990 AUC).
\end{itemize}

The 21\% AUC gain on flow-only data is the key takeaway: raw flow features (ports, protocols, bytes, packets) carry only weak lateral movement signal, but when the same flows are organized into a graph structure, the resulting topology (which hosts connect to which, how densely, and with what temporal patterns) provides near-perfect discrimination. This confirms that lateral movement is fundamentally a relational phenomenon, not a property of individual network flows.

\subsection{False Positive Analysis}
\label{sec:fp_analysis}

Despite high AUC, all methods produce false positives at operationally relevant thresholds. Understanding \emph{why} these false positives occur is critical for deploying lateral movement detection in production. Table~\ref{tab:precision} reports precision values across all methods and variants.

\begin{table}[h]
\centering
\begin{tabular}{lccccc}
\toprule
Variant & Method & AUC & Recall & FPR & Precision \\
\midrule
\multirow{3}{*}{Auth Only}
  & Graph-based      & \textbf{0.984} & 0.922 & 0.006 & \textbf{0.098} \\
  & Isolation Forest & 0.973 & 0.945 & 0.029 & 0.071 \\
  & One-Class SVM    & 0.923 & 0.942 & 0.107 & 0.020 \\
\midrule
\multirow{3}{*}{Combined}
  & Graph-based      & 0.959 & 0.211 & 0.004 & 0.034 \\
  & One-Class SVM    & \textbf{0.800} & 0.052 & \textbf{0.001} & \textbf{0.087} \\
  & Isolation Forest & 0.957 & 0.588 & 0.029 & 0.033 \\
\midrule
\multirow{3}{*}{Flow Only}
  & Graph-based      & \textbf{0.963} & 0.068 & 0.030 & 0.005 \\
  & One-Class SVM    & 0.662 & 0.029 & 0.102 & 0.001 \\
  & Isolation Forest & 0.718 & 0.042 & 0.100 & 0.002 \\
\bottomrule
\end{tabular}
\caption{Precision values for all methods. Low precision reflects extreme class imbalance: hundreds of thousands of benign edges produce false alarms even at low FPR. Graph-based FPR is reported in pair-space; baseline FPR in edge-space.}
\label{tab:precision}
\end{table}

Before examining these specific sources, it helps to ask why topology-based detection produces false positives at all. Lateral movement is, by definition, lateral access: an attacker pivoting through a network uses the same protocols (NTLM, SMB, RDP), the same credentials, and the same network paths as an administrator performing routine maintenance. What separates attack from administration is intent, and intent is not recorded in authentication or flow logs. Our graph features capture the topology of lateral movement (high fan-out from a single source, NTLM-based authentication chains, connections into peripheral hosts), yet legitimate lateral access produces the same topology. The false positives analyzed below are therefore not a defect of our feature design but a property of any detector that reasons from connectivity alone. Administrative activity, service-to-service communication, and machine account authentication each represent a distinct way that benign work can produce attack-shaped topology.

The false positives in our graph-based approach arise from three structural sources:

\textbf{Administrative activity.} System administrators performing legitimate lateral connections (e.g., deploying patches, running remote diagnostics, managing services across multiple hosts) exhibit high fan-out, NTLM usage, and network logon patterns that are indistinguishable from lateral movement at the edge level. In the LANL dataset, these admin patterns account for the majority of false positives in the auth-only variant. The graph structure captures the same connectivity patterns that make lateral movement detectable, but legitimate admin activity shares those same patterns.

\textbf{Service-to-service communication.} In the combined and flow-only variants, automated service-to-service connections (e.g., backup jobs, monitoring agents, load balancer health checks) create dense subgraphs with high edge weights and regular temporal patterns. These mimic the repetitive, high-frequency connection patterns that graph features associate with reconnaissance activity. The flow-only graph, being denser than the auth-only graph, is particularly susceptible to these false positives.

\textbf{Machine account authentication.} Service accounts and machine accounts (identified by the \texttt{\$} naming convention) authenticate to multiple hosts as part of normal domain operations. These create high-degree nodes that the graph model flags as suspicious based on fan-out and degree features. In the LANL dataset, machine accounts are treated as regular nodes, and their normal authentication patterns generate a significant portion of the false positive burden.

These observations suggest several practical mitigation strategies: (1) maintain a whitelist of known administrative source IPs and service accounts to suppress their edges from scoring; (2) incorporate temporal baselines to distinguish periodic service-to-service communication from bursty reconnaissance; and (3) use graph community detection to identify and exclude normal subgraphs (e.g., monitoring infrastructure) from the anomaly scoring pool.

\section{Discussion and Limitations}

Our results show that authentication log data alone provides the strongest signal for lateral movement detection on the LANL-2015 dataset, with the auth-only graph method achieving 92.2\% recall at 0.64\% FPR and 0.984 AUC. Graph-based scoring consistently outperforms unsupervised baselines on all three data variants when applied to the same edge features. The gap is most pronounced on flow-only data, where graph structure provides critical discriminative signal that tabular features alone cannot capture.

Graph-derived features drive most of the detection performance (0.987 AUC vs. 0.953 for tabular-only on the combined variant; 0.998 vs. 0.823 on flow-only). The feature audit confirms that the strongest individual signals are authentication semantics (NTLM usage, network logon type) combined with graph topology (in-degree, out-degree, fan-out).

\textbf{Practical implications.} For a security operations center (SOC) that wants to deploy lateral movement detection, our results point to a concrete configuration. Authentication logs are the primary signal source. They are smaller than flow data and cheaper to ingest, and on this network they were sufficient on their own. Each edge should be scored with graph-derived features rather than treating logins in isolation: fan-out, in-degree, and out-degree capture the spread pattern of a compromised host reaching laterally into new machines, while a temporal burst feature flags hosts that suddenly connect to many new neighbors inside a short window. Edges that sit on multi-hop chains should receive a path-level boost, because lateral movement is a sequence of hops rather than a single jump. In the auth-only configuration this setup reaches the recall and FPR reported above while processing on the order of tens of thousands of authentication events per second, fast enough for near-real-time scoring on a mid-size enterprise network. The main deployment risk is false positives from administrative and service accounts: our false positive analysis (Section~\ref{sec:fp_analysis}) shows that admin hosts and machine accounts produce the same fan-out and NTLM topology as actual compromise, so an unfiltered deployment would spend most of its alerts on routine maintenance. Two mitigations follow from that analysis. Whitelist known administrative source hosts and service-account paths so their edges are excluded from scoring, and build temporal baselines for service traffic so that periodic backup and monitoring jobs stop resembling reconnaissance bursts. With those controls in place, the remaining alert volume stays low enough for a SOC analyst to triage within a normal shift.

Several limitations remain. The LANL dataset represents a single enterprise network over 58 days with red-team attacks that may not capture all real-world lateral movement techniques; 93.6\% of attack events originate from a single compromised host (C17693). Real-world lateral movement may involve more diverse pivot patterns, living-off-the-land techniques, and encrypted channels that do not produce the authentication events our method relies on. The combined method's lower recall relative to auth-only (21.1\% vs. 92.2\%) suggests that our current scoring function does not effectively integrate flow-level signals with authentication semantics. False positives remain a challenge at production scale, though our analysis identifies specific structural causes and suggests targeted mitigations.

\section{Conclusion}

We presented a graph-based approach for lateral movement detection that integrates network flow logs and authentication logs into a unified analysis framework, evaluated on the LANL-2015 dataset. Our key findings are:

\begin{enumerate}
\item \textbf{Auth-only graph analysis achieves the best lateral movement detection:} 92.2\% recall at 0.64\% FPR with 0.984 AUC, outperforming Isolation Forest (0.973 AUC) and One-Class SVM (0.923 AUC) on the same features.

\item \textbf{Graph-based lateral movement detection generalizes across data variants:} it maintains AUC $\geq$ 0.96 on auth-only, combined, and flow-only data, while unsupervised baselines degrade to 0.66--0.72 AUC on flow-only. This demonstrates that modeling lateral movement as a graph problem (reasoning about which hosts connect to which, and how) provides robust signal regardless of which log source is available.

\item \textbf{Graph topology is the primary driver of lateral movement detection:} graph-derived features alone achieve up to 0.998 AUC compared to 0.823 for pure tabular features on flow data, a 21\% gain. This confirms that lateral movement is fundamentally a relational phenomenon: the pattern of connections between hosts carries more discriminative signal than the properties of individual flows or authentication events.

\item \textbf{Lateral movement detection weights generalize to unseen data:} Nelder-Mead optimized weights achieve held-out AUC within 2.1\% of calibration AUC across all variants, matching logistic regression performance on the same features. This confirms that the features themselves carry the lateral movement signal, not the optimization method.

\item \textbf{False positives in lateral movement detection have identifiable structural causes:} administrative activity, service-to-service communication, and machine account authentication produce patterns that mimic lateral movement at the edge level. Targeted mitigations (admin whitelisting, temporal baselining, and community-based exclusion) can reduce the false positive burden without sacrificing detection rate.
\end{enumerate}

The combined graph method's lower recall relative to auth-only suggests that integrating flow and authentication signals requires more sophisticated feature engineering or scoring approaches. Future work should test on additional real-world cloud VPC datasets, explore incremental graph updates for real-time lateral movement detection, and develop scoring functions that better integrate flow-level signals with authentication semantics while accounting for the different noise levels across log sources.

{\footnotesize \bibliographystyle{acm}
\bibliography{sample}}

\end{document}
